\section{Branch analysis and genus}\label{sec:genus}

Define
\begin{equation}\label{eq:q6}
\begin{split}
Q_6(\sigma)={}&64\sigma^6-448\sigma^5+848\sigma^4+80\sigma^3\\
&-1048\sigma^2+152\sigma-151.
\end{split}
\end{equation}

\begin{theorem}\label{thm:genus}
The adopted curve $P(\sigma,x)=0$ is geometrically integral.  Its smooth
projective normalization $\Ccurve$ has genus three.  The map
$\Ccurve\to\Pone_\sigma$ has total ramification 18.
\end{theorem}

\begin{proof}
Specializing $\sigma=-3$ and reducing modulo 2 gives
$x^7+x^6+x^5+x^4+1$, irreducible over $\F_2$.  Since $P$ is monic in $x$,
the generic cover is irreducible over $\Q(\sigma)$.  The nontrivial inertia
below lies in geometric monodromy.  In prime degree seven, the orbits of the
normal geometric subgroup are blocks; nontrivial inertia forces one block,
so the curve is geometrically integral.

Exact elimination gives
\begin{equation}\label{eq:disc}
 \Disc_xP=(4\sigma-9)^2Q_6(\sigma)^3,
 \qquad \Disc Q_6=2^{63}\cdot97.
\end{equation}
The polynomial $Q_6$ is irreducible.  Over
$\Q[\sigma]/(Q_6)$, the gcd of $P$ and $P_x$ has degree three and has no
common root with either $P_{xx}$ or $P_\sigma$.  Consequently each of the six
$Q_6$-values has three distinct smooth points of ramification index two.

At $\sigma=9/4$,
\[
 P(x,9/4)=\frac{(4x-3)(4x-1)^2(16x^2+7)(16x^2-16x+11)}{16384}.
\]
Writing $t=\sigma-9/4$ and $y=x-1/4$, the tangent cone at the unique
singular point is
\[
 -\frac18(y^2-10ty+137t^2),
\]
with discriminant $-7t^2$.  Thus the point is an ordinary node with two
nonvertical tangents; both normalization branches are unramified over $t$.

For infinity put $t=1/\sigma$ and $y=x/\sigma$.  The integral equation
$F(t,y)=t^7P(1/t,y/t)$ has fibre
$(y-1)^4(y+1)^3$ at $t=0$.  Successive blow-up initial forms listed in
Appendix~\ref{app:infinity} separate four branches above $y=1$ and three
above $y=-1$.  Each branch uses $t$ as a uniformizer and is unramified.
Hence the total ramification is $6\cdot3=18$, and Riemann--Hurwitz yields
\[
 2g-2=7(-2)+18=4.
\]
Therefore $g=3$.
\end{proof}

The plane closure is a septic of arithmetic genus 15.  The finite node has
$\delta=1$.  The $t$-adic order 22 of $\Disc_yF$, together with zero
ramification at infinity, gives total infinity defect 11.  Thus
$15-1-11=3$, independently checking Theorem~\ref{thm:genus}.
